%% ch04-signals.tex — Signal Reference
%% JA4proxy Technical Reference Manual

\chapter{Signal Reference}
\index{signals}
\label{ch:signals}

\begin{chapterquote}
A signal is a scored observation. Each of the eight signal checks contributes a score
in [0, 100] and a confidence weight. The risk scorer combines them. No single signal
can block a connection — only the combined score, acting through the dial, produces
blocking actions.
\end{chapterquote}

\newthought{All eight signals run concurrently} after the fast-path bypass zone.
Each signal is independent. An unreachable external service contributes zero — it
never increases the score (fail open). Signal results are logged in the
\code{signals} field of every connection event.

\section{Signal Score Model}
\index{signals!score model}
\index{RiskSignal}

Each signal produces a \code{RiskSignal} with three fields:

\begin{description}
  \item[\code{score}] Integer 0–100. Higher means more suspicious.
  \item[\code{weight}] Float 0.0–1.0. The confidence weight of this signal type.
        High-confidence signals (TLS version, JA4 blacklist category) have higher weights.
        Low-certainty signals (RDAP org, passive DNS) have lower weights.
  \item[\code{reason}] Human-readable explanation for the log.
\end{description}

The risk scorer computes: $\text{final\_score} = \text{clamp}(\sum_i \text{score}_i \times \text{weight}_i, 0, 100)$
where weights are normalised to sum to 1.0 across all active signals.

\begin{fullwidth}
\begin{table}[h]
\caption{Signal summary — all eight signals}
\label{tab:signal-summary}
\begin{tabular}{@{}lllllp{4cm}@{}}
\toprule
\textbf{\#} & \textbf{Signal} & \textbf{Phase} & \textbf{Default weight} & \textbf{External service} & \textbf{Failure behaviour} \\
\midrule
1 & TLS version \& cipher & 3 & 0.20 & None & N/A — local check \\
2 & SNI analysis & 4 & 0.15 & None & N/A — local check \\
3 & TCP \& connection behaviour & 5 & 0.15 & Redis & Score=0 on Redis error \\
4 & ASN classification & 6 & 0.15 & MaxMind DB (mmap) & Score=0 if DB unavailable \\
5 & FCrDNS enrichment & 7 & 0.10 & DNS resolver & Score=0 on DNS error \\
6 & Beaconing detector & 9 & 0.10 & Redis & Score=0 on Redis error \\
7 & AbuseIPDB score & 10 & 0.10 & AbuseIPDB API & Score=0 on API error \\
8 & RDAP org reputation & 11 & 0.05 & RDAP/WHOIS APIs & Score=0 on API error \\
\bottomrule
\end{tabular}
\end{table}
\end{fullwidth}

\section{Signal 1 — TLS Version and Cipher Check}
\index{signals!TLS version}
\index{TLS!version enforcement}
\index{cipher suites}
\label{sig:tls}

\textbf{Phase:} 3 \quad \textbf{Module:} \code{src/security/tls\_enforcer.py}

\subsection{What It Measures}

This signal checks the TLS protocol version and cipher suite list from the ClientHello
against a set of hard rules. TLS 1.0 and 1.1 are deprecated and have known
vulnerabilities; SSLv3 is critically broken. Weak cipher suites indicate either an
outdated client or deliberate downgrade probing.

\subsection{Score Computation}
\index{TLS!SSLv3}
\index{TLS 1.0}
\index{TLS 1.1}

When \configkey{security\_policy.tls\_version\_bypass.enabled} is \code{true}, TLS
1.0/1.1/SSLv3 trigger an immediate TCP RST (before scoring). When the bypass is disabled,
they contribute a RiskSignal instead:

\begin{fullwidth}
\begin{table}[h]
\caption{TLS version score contributions}
\label{tab:tls-scores}
\begin{tabular}{@{}llp{7cm}@{}}
\toprule
\textbf{Condition} & \textbf{Score} & \textbf{Notes} \\
\midrule
TLS 1.3, strong ciphers & 0 & Clean \\
TLS 1.2, ECDHE + AES-GCM / ChaCha20 & 0 & Acceptable \\
TLS 1.2 with weak ciphers (RC4, DES, 3DES, NULL, EXPORT) & 40 & Routed through scorer \\
TLS 1.1 (bypass disabled) & 70 & Deprecated protocol \\
TLS 1.0 (bypass disabled) & 80 & Deprecated protocol \\
SSLv3 (bypass disabled) & 95 & Critically deprecated \\
\bottomrule
\end{tabular}
\end{table}
\end{fullwidth}

\subsection{Configuration Keys}

\configkey{security\_policy.tls\_version\_bypass.enabled} — when true, old TLS versions
trigger TCP RST directly. When false, they contribute a risk score.

\subsection{Failure Behaviour}

This signal performs only local parsing of the ClientHello. There is no external service
dependency. It cannot fail.

\section{Signal 2 — SNI Analysis}
\index{signals!SNI}
\index{SNI}
\index{DGA}
\label{sig:sni}

\textbf{Phase:} 4 \quad \textbf{Module:} \code{src/security/sni\_analyzer.py}

\subsection{What It Measures}

Server Name Indication (SNI) is an extension in the ClientHello that specifies which
hostname the client is trying to reach. Legitimate browsers always include SNI when
connecting to HTTPS services. Missing SNI, IP-literal SNI, and algorithmically generated
domain names are all strong signals of automated tools or malware.

\subsection{Score Computation}

\begin{fullwidth}
\begin{table}[h]
\caption{SNI signal score contributions}
\label{tab:sni-scores}
\begin{tabular}{@{}llp{6cm}@{}}
\toprule
\textbf{Condition} & \textbf{Score} & \textbf{Notes} \\
\midrule
SNI present, matches expected hostname & 0 & Normal browser or API client \\
SNI present, unexpected hostname & 10--30 & May be misconfigured client or proxy \\
SNI absent & 50 & Scanners and tools typically omit SNI \\
SNI is an IP literal (e.g., \code{192.168.1.1}) & 60 & Invalid SNI; tools and scanners \\
SNI is DGA-like (high entropy, no real TLD) & 40--80 & Score based on entropy measurement \\
\bottomrule
\end{tabular}
\end{table}
\end{fullwidth}

\subsection{DGA Detection Algorithm}
\index{DGA!detection algorithm}

Domain Generation Algorithm detection uses Shannon entropy of the domain label
characters. A real domain like \code{example.com} has low entropy; a DGA domain like
\code{xk9p2mq4vrl.cn} has high entropy. The threshold is tunable:
\begin{itemize}
  \item Entropy $<$ 3.5 bits/character: clean
  \item Entropy 3.5–4.2 bits/character: moderate score (40)
  \item Entropy $>$ 4.2 bits/character: high score (80)
\end{itemize}

Additionally, the TLD is checked against a list of TLDs commonly used in DGA campaigns.
A domain matching a flagged TLD with high entropy scores at the upper end of the range.

\subsection{Failure Behaviour}

This signal performs only local computation. There is no external service dependency.
It cannot fail.

\section{Signal 3 — TCP and Connection Behaviour}
\index{signals!TCP behaviour}
\index{JA4T}
\index{session resumption}
\index{concurrency}
\label{sig:tcp}

\textbf{Phase:} 5 \quad \textbf{Modules:} \code{src/security/tcp\_analyzer.py}, \code{src/security/mtls.py}

\subsection{What It Measures}

TCP and connection behaviour signals capture how a client connects over time, not just
what it sends. Automated tools have different connection patterns from real users:
they connect more frequently, resume sessions less often, and maintain more simultaneous
connections.

\subsection{JA4T Fingerprinting}
\index{JA4T!definition}

JA4T is the TLS transport fingerprint — a hash of the TCP window size, TCP options
(MSS, window scale, SACK, timestamps), and their ordering. Unlike JA4 (which fingerprints
the TLS application layer), JA4T fingerprints the TCP/IP stack implementation. Together,
JA4 + JA4T can distinguish a real Chrome browser from a tool that has copied Chrome's
JA4 value but runs on a different OS or with a modified TCP stack.

\subsection{Score Components}

\begin{fullwidth}
\begin{table}[h]
\caption{TCP behaviour signal score components}
\label{tab:tcp-scores}
\begin{tabular}{@{}llp{6.5cm}@{}}
\toprule
\textbf{Component} & \textbf{Score contribution} & \textbf{Notes} \\
\midrule
JA4T mismatch with claimed browser & +20 & JA4T profile inconsistent with JA4 fingerprint label \\
Session resumption ratio $<$ 10\% & +15 & Legitimate browsers resume sessions aggressively \\
Concurrent connections $>$ 20 & +30 & Legitimate browsers maintain 6–8 connections \\
Very short connection lifespan ($<$ 0.1s) & +20 & Scanner pattern — connect, sniff, disconnect \\
First-seen IP (no return visit history) & +5 & Weak signal; most IPs start here \\
Return visitor with all blocked history & +25 & IP that only ever got blocked \\
\bottomrule
\end{tabular}
\end{table}
\end{fullwidth}

\subsection{Redis Keys Used}

\begin{itemize}
  \item \code{sess:\{ip\}} — Hash of \code{total} and \code{resumed} session counts
  \item \code{conn:\{ip\}} — Current concurrent connection count (INCR/DECR)
  \item \code{visit:\{ip\}} — Hash of \code{first\_seen}, \code{total}, \code{allowed}, \code{blocked}
\end{itemize}

\subsection{Failure Behaviour}

If Redis is unreachable, all Redis-backed sub-signals return score=0. The signal as a
whole returns score=0. The proxy logs the Redis error and increments
\code{ja4proxy\_redis\_errors\_total}.

\section{Signal 4 — ASN Classification}
\index{signals!ASN}
\index{ASN}
\index{datacenter detection}
\index{Tor!exit node}
\label{sig:asn}

\textbf{Phase:} 6 \quad \textbf{Module:} \code{src/security/asn\_classifier.py}

\subsection{What It Measures}

The Autonomous System Number (ASN) associated with the client IP identifies the network
organisation that owns the IP block. Legitimate browser traffic overwhelmingly originates
from residential ISPs, mobile carriers, and corporate networks. Automated scanners and
bots are disproportionately hosted in datacenters and cloud providers.

\subsection{Score Computation}

ASN classification uses:
\begin{enumerate}
  \item The MaxMind GeoLite2 ASN database (memory-mapped for zero-copy lookup)
  \item A curated datacenter ASN list in \code{config/asn\_datacenter\_list.yml}
  \item The Tor exit node list (fetched periodically from \code{check.torproject.org})
\end{enumerate}

\begin{fullwidth}
\begin{table}[h]
\caption{ASN signal score contributions}
\label{tab:asn-scores}
\begin{tabular}{@{}llp{6cm}@{}}
\toprule
\textbf{ASN category} & \textbf{Score} & \textbf{Notes} \\
\midrule
Residential ISP & 0 & Expected for browser traffic \\
Mobile carrier & 0 & Expected for mobile browser traffic \\
Corporate network & 5 & Minor flag — APIs often originate here \\
Hosting / datacenter (general) & 25 & Common for cloud-hosted scanners \\
Known VPN provider & 30 & Often used to evade geo-blocking \\
Tor exit node & 65 & Significant signal — anonymisation tool \\
Known botnet/abuse ASN & 80 & From curated bad-org list \\
\bottomrule
\end{tabular}
\end{table}
\end{fullwidth}

\subsection{Failure Behaviour}

If the MaxMind database file is missing or corrupt, this signal returns score=0.
The proxy logs the failure and increments \code{ja4proxy\_signal\_duration\_seconds}
with an error tag. The GeoLite2 database should be updated monthly and validated
at startup.

\section{Signal 5 — FCrDNS Enrichment}
\index{signals!FCrDNS}
\index{FCrDNS}
\index{PTR record}
\index{passive DNS}
\label{sig:dns}

\textbf{Phase:} 7 \quad \textbf{Module:} \code{src/security/dns\_enrichment.py}

\subsection{What It Measures}

Forward-confirmed reverse DNS (FCrDNS) validates that the PTR record for a client IP
resolves to a hostname, and that hostname's A/AAAA record resolves back to the same IP.
Residential ISPs, cloud providers, and CDNs all maintain consistent FCrDNS. Ephemeral
scanner infrastructure typically does not.

\subsection{Classification Logic}

\begin{fullwidth}
\begin{table}[h]
\caption{FCrDNS signal classification and scores}
\label{tab:dns-scores}
\begin{tabular}{@{}llp{7cm}@{}}
\toprule
\textbf{Classification} & \textbf{Score} & \textbf{Notes} \\
\midrule
FCrDNS valid, residential pattern & 0 & PTR resolves to ISP hostname pattern \\
FCrDNS valid, cloud/datacenter & 10 & e.g.\ \code{ec2-*.compute.amazonaws.com} \\
PTR absent or NXDOMAIN & 30 & No reverse DNS — common for scanners \\
PTR present but FCrDNS fails & 35 & Inconsistent DNS — suspicious \\
PTR matches known scanning hostname & 60 & e.g.\ Shodan, Censys, known scanner infra \\
\bottomrule
\end{tabular}
\end{table}
\end{fullwidth}

\subsection{Async Lookup Architecture}

DNS lookups are fire-and-forget: \code{asyncio.create\_task(dns\_lookup(ip))} is called
at pipeline start. Results are written to a short-lived in-process cache. If the lookup
has not completed when the scorer runs, the signal returns score=0 for this connection
(fail open). Results are available for subsequent connections from the same IP.

\subsection{Failure Behaviour}

DNS resolution errors (NXDOMAIN, SERVFAIL, timeout) cause this signal to return score=0.
The failure is logged at DEBUG level (DNS failures are common and non-actionable).

\section{Signal 6 — Beaconing Detector}
\index{signals!beaconing}
\index{beaconing}
\index{IAT}
\label{sig:beacon}

\textbf{Phase:} 9 \quad \textbf{Module:} \code{src/security/beaconing\_detector.py}

\subsection{What It Measures}

Beaconing is the pattern of regular, low-jitter connections from the same source —
characteristic of C2 implants and automated scheduled tasks rather than human browsing.
Human browsing has highly variable inter-arrival times (IAT). Beaconing tools produce
IAT distributions with low coefficient of variation.

\subsection{IAT Algorithm}
\index{IAT!coefficient of variation}

\begin{enumerate}
  \item Connection timestamps are stored in Redis Sorted Set \code{beacon:\{ip\}:\{ja4\}}
        (score = Unix timestamp in milliseconds)
  \item The sliding window retains the last $N$ timestamps within a 30-minute window
  \item \textbf{Inter-arrival times} are computed as consecutive differences:
        $\text{IAT}_i = t_{i+1} - t_i$
  \item The \textbf{coefficient of variation} (CV) is computed: $\text{CV} = \sigma / \mu$
  \item Low CV indicates regular beaconing; high CV indicates human browsing
\end{enumerate}

\begin{fullwidth}
\begin{table}[h]
\caption{Beaconing signal score thresholds}
\label{tab:beacon-scores}
\begin{tabular}{@{}llp{6cm}@{}}
\toprule
\textbf{CV range} & \textbf{Score} & \textbf{Notes} \\
\midrule
$<$ 5 connections in window & 0 & Insufficient data; no score assigned \\
CV $>$ 0.8 & 0 & High variability — human pattern \\
CV 0.4–0.8 & 20 & Moderate regularity \\
CV 0.1–0.4 & 50 & Strong beaconing pattern \\
CV $<$ 0.1 & 75 & Near-perfect regularity — almost certainly automated \\
\bottomrule
\end{tabular}
\end{table}
\end{fullwidth}

\subsection{Redis Keys Used}

\code{beacon:\{ip\}:\{ja4\}} — Sorted Set where member = timestamp and score = timestamp.
ZRANGEBYSCORE is used to retrieve the sliding window. ZADD appends new timestamps.
Old members (outside the window) are removed with ZREMRANGEBYSCORE.

\subsection{Failure Behaviour}

If Redis is unreachable, this signal returns score=0. Redis writes are fire-and-forget
and do not block the pipeline.

\section{Signal 7 — AbuseIPDB Score}
\index{signals!AbuseIPDB}
\index{AbuseIPDB}
\label{sig:abuseipdb}

\textbf{Phase:} 10 \quad \textbf{Module:} \code{src/security/abuseipdb.py}

\subsection{What It Measures}

AbuseIPDB is a community-maintained database of IP addresses reported for malicious
activity (scanning, brute-force, spam, C2 traffic). Each IP has a confidence score
(0–100) reflecting how frequently it has been reported and how recently.

\subsection{Score Computation}

\begin{fullwidth}
\begin{table}[h]
\caption{AbuseIPDB signal score mapping}
\label{tab:abuseipdb-scores}
\begin{tabular}{@{}llp{6.5cm}@{}}
\toprule
\textbf{AbuseIPDB confidence} & \textbf{Risk signal score} & \textbf{Notes} \\
\midrule
0–49 & 0 & Below minimum threshold; no signal \\
50–74 & 30 & Moderate abuse history \\
75–89 & 60 & Significant abuse history \\
90–100 & 80 & Very high confidence abuse \\
API error / timeout & 0 & Fail open — never increases score \\
\bottomrule
\end{tabular}
\end{table}
\end{fullwidth}

\subsection{Caching Architecture}
\index{AbuseIPDB!caching}

AbuseIPDB responses are cached in Redis as \code{abuseipdb:\{ip\}} (JSON, TTL=3600s by
default). Cache hits return immediately with zero API call. Cache misses trigger an
async API call (fire-and-forget) and return score=0 for the current connection. The
result is available for subsequent connections from the same IP.

This design means the \emph{first} connection from a never-seen-before IP pays no
AbuseIPDB penalty — it is allowed while the lookup runs. This is intentional (fail open).

\subsection{Failure Behaviour}

API errors, HTTP 429 rate limit responses, and network timeouts all produce score=0.
The error is logged with context and the \code{ja4proxy\_abuseipdb\_lookups\_total\{result="error"\}}
counter is incremented. The proxy continues functioning with zero AbuseIPDB signal.

\section{Signal 8 — RDAP Org Reputation}
\index{signals!RDAP}
\index{RDAP}
\index{WHOIS}
\label{sig:rdap}

\textbf{Phase:} 11 \quad \textbf{Module:} \code{src/security/rdap\_enrichment.py}

\subsection{What It Measures}

RDAP (Registration Data Access Protocol) and WHOIS provide the organisation that owns
an IP block. \ja\ maintains a curated list of known-bad organisations in
\code{config/known\_bad\_orgs.yml}. Matching the RDAP organisation name against this
list provides a signal about whether the client IP is associated with organisations
that systematically operate malicious infrastructure.

\subsection{Score Computation}

\begin{fullwidth}
\begin{table}[h]
\caption{RDAP signal score contributions}
\label{tab:rdap-scores}
\begin{tabular}{@{}llp{6cm}@{}}
\toprule
\textbf{Condition} & \textbf{Score} & \textbf{Notes} \\
\midrule
Known-good org (major cloud, CDN) & 0 & AWS, GCP, Azure, Cloudflare, Akamai \\
Unknown org & 5 & Default for unrecognised registrants \\
Org name matches known-bad list & 60--80 & Score from \code{known\_bad\_orgs.yml} \\
RDAP lookup failed & 0 & Fail open \\
\bottomrule
\end{tabular}
\end{table}
\end{fullwidth}

\subsection{Block Expansion}
\index{RDAP!block expansion}

When \configkey{rdap.block\_expansion: true} (disabled by default), blocking an IP
from a known-bad org triggers an automatic CIDR block of the entire /24 (IPv4) or
/48 (IPv6) prefix containing that IP. The expanded block is written to Redis and picked
up by the Layer 0d CIDR trie refresh.

\begin{dangerbox}[Block Expansion Risk]
Block expansion must never be enabled without manual review of the affected prefix.
RDAP data is sometimes stale, incorrect, or shared between legitimate and malicious
operators. A rogue /24 expansion affecting a major ISP block can lock out thousands
of legitimate users.
\end{dangerbox}

\subsection{WHOIS Pivot}

The RDAP enrichment module also performs WHOIS pivoting: it looks up the registrant
email domain for the IP block's registrant record. A registrant email from a known
throwaway or bulletproof hosting domain adds an additional score contribution.

\subsection{Failure Behaviour}

RDAP API errors, timeouts, and malformed responses all return score=0. RDAP queries
are cached in Redis (\code{rdap:\{ip\_or\_prefix\}}, TTL=86400s) to avoid repeated
lookups for the same prefix.
